% ============================================================
% CAPITOLO 1
% Digital Government, innovazione e cybersecurity nella PA
% ============================================================
\chapter{Digital Government, innovazione e cybersecurity nella PA}
\label{chap:digital-government}

% --- Introduzione al capitolo ---
Il presente capitolo ricostruisce l'evoluzione del paradigma di governo digitale, dall'e-Government al Digital Government, analizzando le principali strategie, infrastrutture e politiche che guidano la trasformazione digitale della Pubblica Amministrazione italiana ed europea. Particolare attenzione è dedicata al Piano Triennale ICT, alla cloud strategy, alla cybersecurity e agli indicatori di maturità digitale.


% ============================================================
\section{Evoluzione del paradigma: da e-Government a Digital Government}
\label{sec:evoluzione-paradigma}

% --- Da completare ---
% Contenuti suggeriti:
% - Definizione di e-Government e Digital Government
% - Evoluzione storica: dalla digitalizzazione dei servizi alla trasformazione dei processi
% - Il paradigma "digital by default" e "once only"
% - Ruolo del PNRR nella transizione digitale
% - Confronto tra approccio italiano ed europeo
% - Riferimenti: Commissione Europea, OECD, AgID


% ============================================================
\section{Piano Triennale ICT 2024--2026}
\label{sec:piano-triennale}

% --- Da completare ---
% Contenuti suggeriti:
% - Obiettivi strategici del Piano Triennale 2024-2026
% - Linee di azione: dati, interoperabilità, piattaforme, sicurezza
% - Ruolo di AgID e del Dipartimento per la Trasformazione Digitale
% - Cloud Italia e Polo Strategico Nazionale (PSN)
% - Indicatori di risultato e monitoraggio
% - Riferimenti: AgID, Piano Triennale, Strategia Cloud Italia


% ============================================================
\section{Infrastrutture digitali e cloud strategy nella PA}
\label{sec:infrastrutture-cloud}

% --- Da completare ---
% Contenuti suggeriti:
% - Classificazione dei dati e servizi (ordinari, critici, strategici)
% - Polo Strategico Nazionale: architettura e governance
% - Cloud qualificato: marketplace AgID/ACN
% - Migrazione al cloud: stato dell'arte e criticità
% - On-premise vs cloud nella PA: criteri di scelta
% - Riferimenti: Strategia Cloud Italia, ACN, PSN


% ============================================================
\section{Cybersecurity e sovranità digitale}
\label{sec:cybersecurity}

% --- Da completare ---
% Contenuti suggeriti:
% - Agenzia per la Cybersicurezza Nazionale (ACN)
% - Strategia Nazionale di Cybersicurezza 2022-2026
% - NIS2 e recepimento italiano
% - Perimetro di Sicurezza Nazionale Cibernetica
% - Sovranità digitale e autonomia tecnologica
% - Implicazioni per la scelta on-premise
% - Riferimenti: ACN, NIS2, Perimetro Cibernetico


% ============================================================
\section{Indicatori europei (DESI, Digital Government Index)}
\label{sec:indicatori-europei}

% --- Da completare ---
% Contenuti suggeriti:
% - DESI (Digital Economy and Society Index): struttura e dimensioni
% - Posizionamento dell'Italia nel DESI
% - Digital Government Index (OECD)
% - eGovernment Benchmark (Commissione Europea)
% - Confronto con altri paesi UE
% - Gap e aree di miglioramento per l'Italia
% - Riferimenti: Commissione Europea, OECD, Eurostat
